\begin{center}
\textsc{Gauss} an \textsc{Olbers}.\quad Göttingen, 14.~Junius 1830.
\end{center}

\par\noindent\hspace*{\parindent}\mbox{.\ .\ .\ .\ .\ .}\quad Die Zeit, die mir in diesem Sommer zu eigener Arbeit übrig bleibt, denke ich der Fortsetzung meiner Abhandlung über die biquadratischen Reste zu widmen, damit diese Arbeit, deren erste Anfänge sich schon von 1805 her datiren, ihrer Vollendung näher komme.
Sie wird wenigstens noch zwei ausgedehnte Abhandlungen erfordern; vorerst werde ich es aber bei einer bewenden lassen, und die erste mir dann wieder zu Theil werdende Musse erst wieder einem andern Gegenstande widmen, wahrscheinlich den theoretischen Methoden der höhern Geodäsie.
Es ist seit geraumer Zeit mein Schicksal gewesen, immer solche Arbeiten aufzunehmen, bei denen sich in der Darstellung nicht \textls[120]{schnell} fortschreiten lässt.\quad\mbox{.\ .\ .\ .\ .\ .}
